\documentclass[11pt]{article}
\usepackage[margin=1in]{geometry}
\usepackage{amsmath,amssymb}
\usepackage[T1]{fontenc}
\usepackage{hyperref}
\usepackage{enumitem}
\usepackage{booktabs}

\title{TBYC Dataset Metrics Reference}
\author{}
\date{\today}

\begin{document}
\maketitle

\section{Scope}
This document describes the metrics implemented in \texttt{src/tbyc\_dataset/metrics}: type matching, metadata matching, tag matching, and summary matching. The equations below are written to match the current code behavior, including edge cases such as empty sets and unmatched-type penalties.

\section{Global Notation}
Let an issue be indexed by $i \in \mathcal{I}$ and an artifact type by $t \in \mathcal{T}$.

\begin{itemize}[leftmargin=*]
\item $H_i$: set of human artifact types present in issue $i$.
\item $L_i$: set of LLM artifact types present in issue $i$.
\item $M_i = H_i \cap L_i$: set of matching types in issue $i$.
\item $F_{i,t}^{H}$, $F_{i,t}^{L}$: human and LLM metadata field-name sets for issue $i$, type $t$.
\item $V_{i,t,f}^{H} = (v_1,\dots,v_m)$: human metadata values for issue $i$, type $t$, field $f$.
\item $V_{i,t,f}^{L} = (u_1,\dots,u_n)$: LLM metadata values for issue $i$, type $t$, field $f$.
\item $G_{i,t}^{H}$, $G_{i,t}^{L}$: human and LLM tag sets for issue $i$, type $t$.
\item $S_{i,t}^{H} = (s_1,\dots,s_a)$, $S_{i,t}^{L} = (r_1,\dots,r_b)$: human and LLM summary strings for issue $i$, type $t$.
\item $\operatorname{combine}(x_1,\dots,x_k)$: newline-joined concatenation of non-empty summary strings.
\item $\operatorname{ratio}(x,y)=\begin{cases}1 & y=0 \text{ and } x=0\\0 & y=0 \text{ and } x>0\\x/y & y>0\end{cases}$.
\item $\operatorname{harm}(p,r)=\begin{cases}0 & p+r=0\\ \frac{2pr}{p+r} & p+r>0\end{cases}$.
\end{itemize}

\section{Type Matching}
\subsection{What It Measures}
Type matching measures whether the model predicts the correct artifact \emph{types}, ignoring summaries, tags, and metadata values.

\subsection{Per-Issue Set Metrics}
For issue $i$, define:
\[
\mathrm{TP}_i = |H_i \cap L_i|,\quad
\mathrm{FP}_i = |L_i \setminus H_i|,\quad
\mathrm{FN}_i = |H_i \setminus L_i|,\quad
\mathrm{U}_i = |H_i \cup L_i|.
\]

The reported metrics are:
\[
\mathrm{Precision}_i = \operatorname{ratio}(\mathrm{TP}_i, |L_i|),
\]
\[
\mathrm{Recall}_i = \operatorname{ratio}(\mathrm{TP}_i, |H_i|),
\]
\[
\mathrm{F1}_i = \operatorname{harm}(\mathrm{Precision}_i,\mathrm{Recall}_i),
\]
\[
\mathrm{Jaccard}_i = \operatorname{ratio}(\mathrm{TP}_i,\mathrm{U}_i).
\]

Short description: these scores measure set overlap between human and predicted type inventories for one issue.

\subsection{Macro Average}
If there are $N$ evaluated issues, then:
\[
\mathrm{MacroPrecision} = \frac{1}{N}\sum_{i=1}^{N}\mathrm{Precision}_i,
\quad
\mathrm{MacroRecall} = \frac{1}{N}\sum_{i=1}^{N}\mathrm{Recall}_i,
\]
\[
\mathrm{MacroF1} = \frac{1}{N}\sum_{i=1}^{N}\mathrm{F1}_i,
\quad
\mathrm{MacroJaccard} = \frac{1}{N}\sum_{i=1}^{N}\mathrm{Jaccard}_i.
\]

Short description: this treats each issue equally.

\subsection{Overall Union Metrics}
The code also forms the global unions:
\[
H^{\cup}=\bigcup_{i \in \mathcal{I}} H_i,
\qquad
L^{\cup}=\bigcup_{i \in \mathcal{I}} L_i,
\]
and computes the same set-overlap metrics on $(H^{\cup},L^{\cup})$.

Short description: this asks whether the model covers the repository-level type vocabulary, not whether it gets issue-by-issue presence right.

\subsection{Per-Type Presence Metrics}
For each type $t$ across issues:
\[
\mathrm{TP}_t = |\{i: t \in H_i \land t \in L_i\}|,
\quad
\mathrm{FP}_t = |\{i: t \notin H_i \land t \in L_i\}|,
\]
\[
\mathrm{FN}_t = |\{i: t \in H_i \land t \notin L_i\}|,
\quad
\mathrm{TN}_t = |\{i: t \notin H_i \land t \notin L_i\}|.
\]

Then:
\[
\mathrm{Precision}_t = \operatorname{ratio}(\mathrm{TP}_t,\mathrm{TP}_t+\mathrm{FP}_t),
\]
\[
\mathrm{Recall}_t = \operatorname{ratio}(\mathrm{TP}_t,\mathrm{TP}_t+\mathrm{FN}_t),
\]
\[
\mathrm{F1}_t = \operatorname{harm}(\mathrm{Precision}_t,\mathrm{Recall}_t),
\]
\[
\mathrm{Jaccard}_t = \operatorname{ratio}(\mathrm{TP}_t,\mathrm{TP}_t+\mathrm{FP}_t+\mathrm{FN}_t).
\]

Short description: this measures how reliably one specific type is detected across issues.

\section{Metadata Matching}
\subsection{What It Measures}
Metadata matching evaluates whether metadata values attached to a given type and field agree between human and LLM annotations. Unlike type matching, this metric compares \emph{strings} using exact counts plus soft similarity.

\subsection{Metadata Value Normalization}
Each metadata value is converted to a string before scoring:
\begin{itemize}[leftmargin=*]
\item strings are stripped,
\item numbers and booleans are stringified,
\item lists are joined with spaces,
\item dictionaries are rendered as sorted \texttt{key: value} pairs joined with \texttt{ ; }.
\end{itemize}

Short description: this produces a flat textual representation so all metadata can be compared by phrase-similarity functions.

\subsection{Phrase Normalization}
Before computing similarity, a phrase $x$ is normalized to $\nu(x)$ by:
\begin{itemize}[leftmargin=*]
\item lowercasing,
\item replacing all non-alphanumeric characters with spaces,
\item collapsing repeated whitespace,
\item trimming leading and trailing whitespace.
\end{itemize}

\subsection{Primitive Similarity Functions}
For normalized phrases $a$ and $b$, let $T(a)$ be the set of whitespace-delimited tokens, and let $C_n(a)$ be the set of character $n$-grams of the padded string ``\verb| |\,$a$\,\verb| |''.

\paragraph{Token F1.}
Let $o = |T(a)\cap T(b)|$. Then:
\[
\operatorname{tokPrec}(a,b)=\operatorname{ratio}(o,|T(b)|),\qquad
\operatorname{tokRec}(a,b)=\operatorname{ratio}(o,|T(a)|),
\]
\[
\operatorname{tokenF1}(a,b)=\operatorname{harm}(\operatorname{tokPrec}(a,b),\operatorname{tokRec}(a,b)).
\]
Measures overlap between unique token sets.

\paragraph{Sequence Ratio.}
\[
\operatorname{sequenceRatio}(a,b)=\texttt{SequenceMatcher}(a,b).\texttt{ratio()}.
\]
Measures surface-form similarity of the normalized strings.

\paragraph{Token Jaccard.}
\[
\operatorname{tokenJaccard}(a,b)=\operatorname{ratio}(|T(a)\cap T(b)|,|T(a)\cup T(b)|).
\]
Measures set overlap between token vocabularies.

\paragraph{Character 3-gram Jaccard.}
\[
\operatorname{char3Jaccard}(a,b)=\operatorname{ratio}(|C_3(a)\cap C_3(b)|,|C_3(a)\cup C_3(b)|).
\]
Measures fuzzy overlap at the subword character level.

\paragraph{Token Containment.}
\[
\operatorname{containment}(a,b)=\operatorname{ratio}(|T(a)\cap T(b)|,\min(|T(a)|,|T(b)|)).
\]
Measures whether the shorter token set is mostly contained in the longer one.

\paragraph{Max-All Ensemble.}
\[
\operatorname{maxAll}(a,b)=\max\{
\operatorname{tokenF1}(a,b),
\operatorname{sequenceRatio}(a,b),
\operatorname{tokenJaccard}(a,b),
\operatorname{char3Jaccard}(a,b),
\operatorname{containment}(a,b)
\}.
\]
Measures best-case phrase agreement across several lightweight similarity functions.

\subsection{Greedy One-to-One Matching}
For one field $f$ of one type $t$ in one issue $i$, compute all pairwise similarities:
\[
\sigma_{jk}=\operatorname{sim}(v_j,u_k),
\quad v_j \in V_{i,t,f}^{H},\ u_k \in V_{i,t,f}^{L}.
\]
Sort all candidate pairs by descending $\sigma_{jk}$ and greedily keep a pair only if neither value has been used before. This produces a one-to-one set of greedy pairs $P$.

Short description: this approximates bipartite matching while preferring the strongest alignments first.

\subsection{Hard Metadata Metrics}
Given a threshold $\tau$ (default $0.82$), define:
\[
\mathrm{TP}_{i,t,f}=|\{(j,k)\in P : \sigma_{jk}\ge\tau\}|,
\]
\[
\mathrm{FP}_{i,t,f}=|V_{i,t,f}^{L}|-\mathrm{TP}_{i,t,f},
\qquad
\mathrm{FN}_{i,t,f}=|V_{i,t,f}^{H}|-\mathrm{TP}_{i,t,f}.
\]

Then:
\[
\mathrm{Precision}_{i,t,f}=\operatorname{ratio}(\mathrm{TP}_{i,t,f},|V_{i,t,f}^{L}|),
\]
\[
\mathrm{Recall}_{i,t,f}=\operatorname{ratio}(\mathrm{TP}_{i,t,f},|V_{i,t,f}^{H}|),
\]
\[
\mathrm{F1}_{i,t,f}=\operatorname{harm}(\mathrm{Precision}_{i,t,f},\mathrm{Recall}_{i,t,f}),
\]
\[
\mathrm{Jaccard}_{i,t,f}=\operatorname{ratio}(\mathrm{TP}_{i,t,f},\mathrm{TP}_{i,t,f}+\mathrm{FP}_{i,t,f}+\mathrm{FN}_{i,t,f}).
\]

Short description: these scores count how many metadata values can be matched strongly enough to count as correct.

\subsection{Soft Metadata Metrics}
For each LLM value, compute its best similarity against any human value:
\[
b^{L}(u_k)=\max_{v_j \in V_{i,t,f}^{H}} \operatorname{sim}(v_j,u_k).
\]
For each human value, compute its best similarity against any LLM value:
\[
b^{H}(v_j)=\max_{u_k \in V_{i,t,f}^{L}} \operatorname{sim}(v_j,u_k).
\]

The code stores the sums:
\[
\Sigma^{L}_{i,t,f}=\sum_{u_k \in V_{i,t,f}^{L}} b^{L}(u_k),
\qquad
\Sigma^{H}_{i,t,f}=\sum_{v_j \in V_{i,t,f}^{H}} b^{H}(v_j).
\]

Then:
\[
\mathrm{SoftPrecision}_{i,t,f}=
\begin{cases}
1 & |V_{i,t,f}^{L}|=0\\
\Sigma^{L}_{i,t,f}/|V_{i,t,f}^{L}| & \text{otherwise}
\end{cases}
\]
\[
\mathrm{SoftRecall}_{i,t,f}=
\begin{cases}
1 & |V_{i,t,f}^{H}|=0\\
\Sigma^{H}_{i,t,f}/|V_{i,t,f}^{H}| & \text{otherwise}
\end{cases}
\]
\[
\mathrm{SoftF1}_{i,t,f}=\operatorname{harm}(\mathrm{SoftPrecision}_{i,t,f},\mathrm{SoftRecall}_{i,t,f}).
\]

Short description: these scores measure closeness even when values do not cross the hard-match threshold.

\subsection{Pooling and Aggregation}
Within an issue or type, the code pools counts and similarity sums by simple addition:
\[
\mathrm{TP}=\sum \mathrm{TP}_{i,t,f},\quad
\mathrm{FP}=\sum \mathrm{FP}_{i,t,f},\quad
\mathrm{FN}=\sum \mathrm{FN}_{i,t,f},
\]
\[
\Sigma^L=\sum \Sigma^L_{i,t,f},\qquad \Sigma^H=\sum \Sigma^H_{i,t,f},
\]
and then recomputes:
\[
\mathrm{Precision}=\operatorname{ratio}(\mathrm{TP},N_L),\quad
\mathrm{Recall}=\operatorname{ratio}(\mathrm{TP},N_H),\quad
\mathrm{F1}=\operatorname{harm}(\mathrm{Precision},\mathrm{Recall}),
\]
\[
\mathrm{Jaccard}=\operatorname{ratio}(\mathrm{TP},\mathrm{TP}+\mathrm{FP}+\mathrm{FN}),
\]
\[
\mathrm{SoftPrecision}=
\begin{cases}
1 & N_L=0\\
\Sigma^L/N_L & \text{otherwise}
\end{cases},
\quad
\mathrm{SoftRecall}=
\begin{cases}
1 & N_H=0\\
\Sigma^H/N_H & \text{otherwise}
\end{cases},
\]
\[
\mathrm{SoftF1}=\operatorname{harm}(\mathrm{SoftPrecision},\mathrm{SoftRecall}).
\]
Here $N_H$ and $N_L$ are pooled human and LLM metadata-value counts.

Short description: pooled scores are micro-style metrics over all values in the aggregation bucket.

\subsection{Metadata Matching Outputs}
\begin{itemize}[leftmargin=*]
\item \texttt{per\_issue}: pooled metadata metrics within each issue.
\item \texttt{overall.pooled}: pooled metadata metrics after unioning all issue values by type and field.
\item \texttt{macro\_average}: arithmetic mean of per-issue pooled metrics.
\item \texttt{metadata\_matching\_all.json}: repeats the full computation for each similarity backend and stores a side-by-side summary.
\end{itemize}

\section{Tag Matching}
\subsection{What It Measures}
Tag matching compares human and LLM tag sets, but only for types that appear on both sides of an issue.

\subsection{Per-Issue, Per-Type Tag Overlap}
For issue $i$ and matched type $t \in M_i$, let:
\[
\mathrm{Int}_{i,t}=|G_{i,t}^{H}\cap G_{i,t}^{L}|,
\qquad
\mathrm{Uni}_{i,t}=|G_{i,t}^{H}\cup G_{i,t}^{L}|.
\]

Then:
\[
\mathrm{Precision}_{i,t}=\operatorname{ratio}(\mathrm{Int}_{i,t},|G_{i,t}^{L}|),
\]
\[
\mathrm{Recall}_{i,t}=\operatorname{ratio}(\mathrm{Int}_{i,t},|G_{i,t}^{H}|),
\]
\[
\mathrm{F1}_{i,t}=\operatorname{harm}(\mathrm{Precision}_{i,t},\mathrm{Recall}_{i,t}),
\]
\[
\mathrm{Jaccard}_{i,t}=\operatorname{ratio}(\mathrm{Int}_{i,t},\mathrm{Uni}_{i,t}).
\]

Short description: this measures whether the detailed labels attached to correctly detected types agree.

\subsection{Pooled Tag Metrics}
Within an issue, per type, or overall, the code sums:
\[
N_H=\sum |G_{i,t}^{H}|,\quad
N_L=\sum |G_{i,t}^{L}|,\quad
\mathrm{Int}=\sum \mathrm{Int}_{i,t},\quad
\mathrm{Uni}=\sum \mathrm{Uni}_{i,t},
\]
then computes:
\[
\mathrm{Precision}=\operatorname{ratio}(\mathrm{Int},N_L),\quad
\mathrm{Recall}=\operatorname{ratio}(\mathrm{Int},N_H),
\]
\[
\mathrm{F1}=\operatorname{harm}(\mathrm{Precision},\mathrm{Recall}),\quad
\mathrm{Jaccard}=\operatorname{ratio}(\mathrm{Int},\mathrm{Uni}).
\]

Short description: this is a pooled overlap score over all compared type-tag sets.

\subsection{Macro Average}
The macro average in \texttt{tag\_matching.json} is the arithmetic mean of the per-issue pooled precision, recall, F1, and Jaccard.

\section{Summary Matching}
\subsection{What It Measures}
Summary matching compares the semantic similarity of aggregated human and LLM summaries, but only for matching types.

\subsection{Summary Construction}
For issue $i$ and type $t \in M_i$:
\[
\bar{S}_{i,t}^{H} = \operatorname{combine}(S_{i,t}^{H}),\qquad
\bar{S}_{i,t}^{L} = \operatorname{combine}(S_{i,t}^{L}).
\]

Short description: all summaries of the same type are concatenated into one reference text and one candidate text before scoring.

\subsection{CodeBERT Cosine}
Let $\phi(x)$ be the mean pooled token embedding produced by the configured CodeBERT encoder after truncation to 512 tokens and masking out padding. Then:
\[
\mathrm{CodeBERTCosine}(x,y)=
\frac{\phi(x)\cdot\phi(y)}{\|\phi(x)\|_2\|\phi(y)\|_2}.
\]
The implementation clips the result to $[-1,1]$ and returns $1$ for two empty strings, $0$ if exactly one side is empty.

Short description: this measures embedding-level semantic similarity of the aggregated summaries.

\subsection{BERTScore}
Let $R=\{r_1,\dots,r_m\}$ be the non-special reference token embeddings and $C=\{c_1,\dots,c_n\}$ the non-special candidate token embeddings after L2 normalization. Form the cosine-similarity matrix:
\[
A_{jk}=c_j^\top r_k.
\]
Then:
\[
\mathrm{BERTScorePrecision}(x,y)=\frac{1}{n}\sum_{j=1}^{n}\max_k A_{jk},
\]
\[
\mathrm{BERTScoreRecall}(x,y)=\frac{1}{m}\sum_{k=1}^{m}\max_j A_{jk},
\]
\[
\mathrm{BERTScoreF1}(x,y)=\operatorname{harm}(\mathrm{BERTScorePrecision}(x,y),\mathrm{BERTScoreRecall}(x,y)).
\]
The implementation returns all ones for two empty strings and all zeros if exactly one side is empty.

Short description: this measures token-level semantic coverage between the generated and reference summaries.

\subsection{BLEURT}
Let $\beta_{\mathrm{raw}}(x,y)$ be the scalar output of the configured BLEURT sequence-classification model. The reported score depends on post-processing mode.

\paragraph{Raw mode.}
\[
\beta(x,y)=\beta_{\mathrm{raw}}(x,y).
\]

\paragraph{Clip mode.}
For clipping floor $c_{\min}$:
\[
\beta(x,y)=
\begin{cases}
\max(c_{\min},\beta_{\mathrm{raw}}(x,y)) & c_{\min}\neq \varnothing\\
\beta_{\mathrm{raw}}(x,y) & \text{otherwise}.
\end{cases}
\]

\paragraph{Sigmoid mode.}
For temperature $T>0$ and bias $b$:
\[
z=\frac{\beta_{\mathrm{raw}}(x,y)-b}{T},
\qquad
z \leftarrow \min(60,\max(-60,z)),
\]
\[
\beta(x,y)=\frac{1}{1+e^{-z}}.
\]
The implementation returns $1$ for two empty strings and $0$ if exactly one side is empty.

Short description: BLEURT is a learned text-quality score, optionally normalized to a bounded range for easier aggregation.

\subsection{Issue-Level Coverage Metrics}
For issue $i$:
\[
h_i = |H_i|,\qquad
\ell_i = |L_i|,\qquad
m_i = |M_i|.
\]

Coverage metrics are:
\[
\mathrm{CovPrecision}_i=\operatorname{ratio}(m_i,\ell_i),\quad
\mathrm{CovRecall}_i=\operatorname{ratio}(m_i,h_i),\quad
\mathrm{CovF1}_i=\operatorname{harm}(\mathrm{CovPrecision}_i,\mathrm{CovRecall}_i).
\]

Short description: these scores measure how much of the type inventory was actually available for summary comparison.

\subsection{Issue-Level Summary Scores}
For issue $i$, let $Q_i$ be the set of score rows for matched types. Each row contains:
\[
q_{i,t} =
\bigl(
\mathrm{CodeBERTCosine}_{i,t},
\mathrm{BERTScorePrecision}_{i,t},
\mathrm{BERTScoreRecall}_{i,t},
\mathrm{BERTScoreF1}_{i,t},
\beta_{i,t}
\bigr).
\]

The issue-level matched-only macro score is the mean over matched types:
\[
\overline{q}_i = \frac{1}{|Q_i|}\sum_{q \in Q_i} q,
\]
with all-zero scores if $|Q_i|=0$.

The issue-level weighted score uses weight
\[
w_{i,t}=\max(1,|S_{i,t}^{H}|),
\]
and computes:
\[
\widetilde{q}_i=\frac{\sum_{t \in M_i} w_{i,t} q_{i,t}}{\sum_{t \in M_i} w_{i,t}},
\]
again defaulting to zeros when no matched types exist.

Short description: the macro version weights each matched type equally, while the weighted version gives more influence to types with more human summary chunks.

\subsection{Penalty for Unmatched Types}
For any score vector $q$ and penalty $p \in [0,1]$, define:
\[
\operatorname{PenaltyApply}(q,p)=p \cdot q.
\]
The code uses:
\[
q^{\mathrm{pen}}_i = \operatorname{PenaltyApply}(\overline{q}_i,\mathrm{CovRecall}_i).
\]

Short description: this discounts summary quality when many human types were missed entirely.

\subsection{Per-Type Summary Scores Across Issues}
For each type $t$, all human summaries from issues where $t$ exists on both sides are concatenated together, and all LLM summaries from those same issues are concatenated together:
\[
\bar{S}_{t}^{H}=\operatorname{combine}\Bigl(\bigcup_{i:\, t \in H_i \cap L_i} S_{i,t}^{H}\Bigr),
\]
\[
\bar{S}_{t}^{L}=\operatorname{combine}\Bigl(\bigcup_{i:\, t \in H_i \cap L_i} S_{i,t}^{L}\Bigr).
\]
These are scored once with the same summary scorers. Type-level coverage is:
\[
\mathrm{TypeCovPrecision}_t=\operatorname{ratio}(m_t,\ell_t),\quad
\mathrm{TypeCovRecall}_t=\operatorname{ratio}(m_t,h_t),\quad
\mathrm{TypeCovF1}_t=\operatorname{harm}(\mathrm{TypeCovPrecision}_t,\mathrm{TypeCovRecall}_t),
\]
where $h_t$ is the number of issues containing type $t$ in human data, $\ell_t$ is the number containing it in LLM data, and $m_t$ is the number containing it in both.

The penalized per-type score is:
\[
q^{\mathrm{pen}}_t = \operatorname{PenaltyApply}(q_t,\mathrm{TypeCovRecall}_t).
\]

\subsection{Overall Summary Aggregates}
The code reports several top-level aggregates in \texttt{summary\_matching.json}.

\paragraph{\texttt{matched\_only\_macro}.}
Arithmetic mean of issue-level matched-only score vectors over issues with $m_i>0$.

\paragraph{\texttt{all\_issues\_macro\_raw}.}
Arithmetic mean of issue-level matched-only score vectors over all issues, including zero rows for issues with no matched types.

\paragraph{\texttt{all\_issues\_macro\_with\_unmatched\_penalty}.}
Arithmetic mean of penalized issue-level score vectors $q^{\mathrm{pen}}_i$ over all issues.

\paragraph{\texttt{support\_weighted\_matched\_only}.}
Weighted average of issue-level weighted score vectors $\widetilde{q}_i$ using weight $\max(0,m_i)=m_i$:
\[
q^{\mathrm{supp}}=\frac{\sum_i m_i \widetilde{q}_i}{\sum_i m_i}.
\]

\paragraph{\texttt{coverage}.}
Macro average of issue-level coverage metrics:
\[
\overline{\mathrm{CovPrecision}}=\frac{1}{N}\sum_i \mathrm{CovPrecision}_i,\quad
\overline{\mathrm{CovRecall}}=\frac{1}{N}\sum_i \mathrm{CovRecall}_i,\quad
\overline{\mathrm{CovF1}}=\frac{1}{N}\sum_i \mathrm{CovF1}_i.
\]

\paragraph{\texttt{penalized\_overall}.}
Apply the macro coverage recall as a global penalty:
\[
q^{\mathrm{pen\_overall}}=
\operatorname{PenaltyApply}\bigl(q^{\mathrm{matched\_only}},\overline{\mathrm{CovRecall}}\bigr).
\]

\paragraph{\texttt{per\_type\_matched\_only\_macro}.}
Arithmetic mean of type-level score vectors over types with at least one matched issue.

\paragraph{\texttt{per\_type\_macro\_with\_unmatched\_penalty}.}
Arithmetic mean of penalized type-level score vectors over all types.

\section{Edge-Case Conventions}
Several metrics use the same conventions:
\begin{itemize}[leftmargin=*]
\item empty-vs-empty set comparisons score as $1$ for precision/recall-style ratios when both numerator and denominator are zero,
\item empty-vs-nonempty scores as $0$ when the denominator is zero but the numerator is positive,
\item harmonic means are defined as $0$ when both inputs are zero,
\item summary scorers return perfect agreement for two empty strings and zero for one empty string versus one non-empty string.
\end{itemize}

These conventions matter because some metric buckets can be empty after filtering by issue, type, or matched-type constraints.

\section{File Mapping}
\begin{center}
\begin{tabular}{ll}
\toprule
Output file & Main contents \\
\midrule
\texttt{type\_matching.json} & type presence overlap metrics \\
\texttt{metadata\_matching.json} & hard and soft metadata-value matching metrics \\
\texttt{metadata\_matching\_all.json} & metadata metrics for all similarity backends \\
\texttt{tag\_matching.json} & tag overlap metrics on matching types \\
\texttt{summary\_matching.json} & semantic summary similarity plus coverage-aware aggregates \\
\bottomrule
\end{tabular}
\end{center}

\section{Visualization-Derived Global Points}
The bar chart \texttt{global\_points\_bar.png} in \texttt{src/tbyc\_dataset/metrics/visualization.py} uses a \emph{derived visualization score}, not a separately stored evaluation metric JSON.

\subsection{Component Set}
For each repository-model bundle, the visualization reads these seven component scores:
\[
\mathcal{C}=
\{
\texttt{metadata\_f1},
\texttt{metadata\_soft\_f1},
\texttt{tag\_f1},
\texttt{type\_f1},
\texttt{summary\_bleurt},
\texttt{summary\_bertscore\_f1},
\texttt{summary\_codebert}
\}.
\]

More explicitly, for repository $r$ and model $m$, the component vector is:
\[
c_{r,m} =
\bigl(
\mathrm{MetaF1}_{r,m},
\mathrm{MetaSoftF1}_{r,m},
\mathrm{TagF1}_{r,m},
\mathrm{TypeF1}_{r,m},
\mathrm{SummBLEURT}_{r,m},
\mathrm{SummBERTF1}_{r,m},
\mathrm{SummCodeBERT}_{r,m}
\bigr),
\]
where:
\begin{itemize}[leftmargin=*]
\item $\mathrm{TypeF1}_{r,m}$ is \texttt{type\_matching.json -> macro\_average.f1},
\item $\mathrm{MetaF1}_{r,m}$ is \texttt{metadata\_matching.json -> macro\_average.f1},
\item $\mathrm{MetaSoftF1}_{r,m}$ is \texttt{metadata\_matching.json -> macro\_average.soft\_f1},
\item $\mathrm{TagF1}_{r,m}$ is \texttt{tag\_matching.json -> overall.f1},
\item $\mathrm{SummBERTF1}_{r,m}$ is \texttt{summary\_matching.json -> overall.all\_issues\_macro\_with\_unmatched\_penalty.bertscore.f1},
\item $\mathrm{SummBLEURT}_{r,m}$ is \texttt{summary\_matching.json -> overall.all\_issues\_macro\_with\_unmatched\_penalty.bleurt.score},
\item $\mathrm{SummCodeBERT}_{r,m}$ is \texttt{summary\_matching.json -> overall.all\_issues\_macro\_with\_unmatched\_penalty.codebert.cosine}.
\end{itemize}

\subsection{Per-Bundle Composite Score}
The repository-model composite score is the unweighted arithmetic mean of the seven components:
\[
\mathrm{BundleScore}_{r,m}=\frac{1}{7}\sum_{k \in \mathcal{C}} c_{r,m,k}.
\]

Short description: each repository contributes one equally weighted composite score per model, averaging type, metadata, tag, and penalized summary quality.

\subsection{Global Model Score}
If a model $m$ appears on repositories $\mathcal{R}_m$, the global score used for the bar chart is:
\[
\mathrm{GlobalScore}_{m}=
\frac{1}{|\mathcal{R}_m|}
\sum_{r \in \mathcal{R}_m}\mathrm{BundleScore}_{r,m}.
\]

Short description: this treats each repository equally, regardless of issue count.

\subsection{Conversion to Discrete Points}
Let $P_{\max}$ be \texttt{points\_max} and $\Delta$ be \texttt{points\_step}. The raw points are:
\[
\mathrm{RawPoints}_m = \mathrm{GlobalScore}_m \cdot P_{\max}.
\]
The displayed point total is quantized to the nearest multiple of $\Delta$:
\[
\mathrm{QuantizedPoints}_m =
\Delta \cdot \operatorname{round}\!\left(\frac{\mathrm{RawPoints}_m}{\Delta}\right).
\]
Finally it is clipped to the interval $[0,P_{\max}]$:
\[
\mathrm{Points}_m=
\min\!\bigl(P_{\max},\max(0,\mathrm{QuantizedPoints}_m)\bigr).
\]

Short description: the chart converts a normalized average score into an integer-style point total for presentation.

\subsection{Defaults}
The visualization entry point uses:
\[
P_{\max}=100,\qquad \Delta=1.
\]
With the defaults, the displayed score is simply the global composite score on a 0--100 scale rounded to the nearest integer.

\end{document}
